%------------------------------------------------------------------------%
%----------------   Interne DSLs aus EBNF-Grammatiken     ---------------%
%------------------------------------------------------------------------%

\documentclass[12pt,a4paper,oneside]{scrreprt}        % "scrreprt" is a more european A4 template for reports
\pdfminorversion=4                                    % tell pdflatex to generate PDF in version 1.4
\pdfpagewidth=\paperwidth%
\pdfpageheight=\paperheight%
\pdfobjcompresslevel=0                                % Prohibit LZW Compression
\pdfinclusioncopyfonts=1                              % Copy fonts

\input{glyphtounicode.tex}                            % Mapping to unicode
\input{glyphtounicode-cmr.tex}
\pdfgentounicode=1


%------------------------      OWN COMMANDS      ------------------------%
\providecommand{\main}{.}                             % Relative path to the main file, to that file pathes in subfiles stay correct
\newcommand*\lstinputpath[1]{\lstset{inputpath=#1}}   % Set path to the souce code listings
\newcommand{\gqq}[1]{\glqq #1\grqq{}}                 % German quotation marks
\newcommand{\gq}[1]{\glq #1\grq{}}


%--------------  Information about the author and company  --------------%
\title{Erstellung interner DSLs aus EBNF-Grammatiken}
\author{Etienne~Gramlich}
\newcommand{\topic}{Erstellung interner DSLs aus EBNF-Grammatiken}
\newcommand{\Subject}{Master Thesis}
\newcommand{\keywords}{HTWG, Java, DSL, Interne DSL}
\newcommand{\documentLanguage}{de}
\newcommand{\authorName}{Etienne Gramlich}
\newcommand{\authorBirthday}{10.07.1993}
\newcommand{\authorBirthdayCity}{Singen}
\newcommand{\closingdate}{01.04.2020}
\newcommand{\issuedate}{1.10.2019}
\newcommand{\authorCity}{Gailingen}
\newcommand{\authorCityNumber}{78262}
\newcommand{\authorStreet}{Auestr. 9b}
\newcommand{\authorFederalState}{Baden-Württemberg}
\newcommand{\authorCountry}{Deutschland}
\newcommand{\authorStudyTopic}{Master Informatik - Software Engineering}
\newcommand{\autorEmail}{etgramlich@gmx.de}
\newcommand{\proffessor}{Prof.\ Dr.\ Heiko Drachenfels}
\newcommand{\university}{HTWG Konstanz}
\newcommand{\universityName}{Hochschule f\"{u}r Technik, Wirtschaft und Gestaltung}
\newcommand{\universityCity}{Konstanz}
\newcommand{\universityCityNumber}{78462}
\newcommand{\universityStreet}{Brauneggerstr. 55}
\newcommand{\universityFaculty}{Fakult\"{a}t Informatik}
\newcommand{\carer}{Marko Boger}
\newcommand{\companyName}{HTWG Konstanz}
\newcommand{\companyCity}{Konstanz}
\newcommand{\companyCityNumber}{78462}
\newcommand{\companyCityStreet}{Brauneggerstr. 55}


%---------------------------     PACKAGES     ---------------------------%
\usepackage{subfiles}                          % So that chapters in included files can be comipled separately
\usepackage[english,ngerman]{babel}            % German syllables, plus english (for URLs)
\usepackage[utf8]{inputenc}	                   % UTF-8 files
\usepackage{lmodern}                           % Include Latin Modern T1 font
\usepackage[T1]{fontenc}                       % Use type 1 fonts
\usepackage{csquotes}                          % Multi-language quotes
\usepackage{subcaption}                        % Subfigure
\usepackage{graphicx}                          % Include images
\graphicspath{{10_Pictures/}{../10_Pictures/}} % Specify image location (for main and subfile)
\usepackage{textcomp}
\usepackage[hypcap=true]{caption}              % Rreferences above tables/images
\usepackage{relsize}                           % Dependency of option "smaller" of "acronym"
\usepackage[printonlyused,smaller]{acronym}    % To print only used acronyms
\usepackage{tabulary}				           % Specify width of table rows
\usepackage{booktabs}                          % Better table rules
\usepackage[natbib,         % Allow natnbib commands
            sorting=nty,    % Sort by Name, Title, year
            style=authoryear]{biblatex}
\addbibresource{literature.bib}
\usepackage{scrhack}                           % Avoids warning that \float@addtolists is deprecated in KOMA
\usepackage{listings}                          % Print source code in the document
\lstinputpath{\main/11Code/}                   % Set directory for code listings
\renewcommand{\lstlistlistingname}{Listingverzeichnis}
\usepackage[usenames]{xcolor}                  % Define and use colors
\usepackage{hyperxmp}                          % XMP data in hypersetup
\PassOptionsToPackage{hyphens}{url}            % Allows hyphenation of URLs
\usepackage[unicode,                % Use unicode
			pdftex,                 % Use pdfTeX driver
			bookmarks,              % Creates bookmarks
			bookmarksnumbered=true, % Numbers in PDF bookmarks
			bookmarksopen=true,     % Don't show bookmaks when opened
			hidelinks,              % Links in black and without borders
			hyperindex=true,        % Numbers in index are links
			pdfpagelabels=true,     % Set PDF page labels
			pdfa]{hyperref}                    % Clickable links and TOC
\usepackage[xindy,toc]{glossaries}             % Glossaries (plus better acronyms)
\makeglossaries%
% Glossaries to be used from 'glossaries' package

\newglossaryentry{ANTLR4}{%
name={ANTLR 4},
description={Ein in Java geschriebener Parsergenerator, der verschiedene Zeilsprachen unterstützt (z.B. Java, C++, C\#, Python). Es werden auch LL(*)-Grammatiken unterstützt}
}

\newglossaryentry{Domaenenspezifische Sprache}{%
name={Dom\"anenspezifische Sprache},
description={Eine Programmiersprache, die für eine spezifische Domäne zugeschnitten wurde und somit nur Probleme der Domäne, und keine darüber hinaus, beschreiben lassen. Dadurch können Domänenexperten einfach damit arbeiten und syntaktischer Ballast (Boilerplate Code) wird verhindert}
}

\newglossaryentry{Hostsprache}{%
name={Hostsprache},
description={Eine General-Purpose Programmiersprache, in die eine interne DSL eingebettet ist}
}

\newglossaryentry{Java Virtual Machine}{%
name={Java Virtual Machine},
description={Die Laufzeitumgebung der Sprache Java; es gibt verschiedene Implementierungen für unterschiedliche Plattformen. Darunter sind Windows, Linux Mac OS X, Android und Solaris}
}

\newglossaryentry{Kontextfreie Grammatik}{%
name={Kontextfreie Grammatik},
description={Kontextfreie Grammatiken (\ac{CFG}) sind identisch mit den Typ-2-Grammatiken. Es folgt auf eine Folge von Terminal- und Nichtterminalsymbolen immer ein Nichtterminalsymbol}
}

\newglossaryentry{Regulaere Grammatik}{%
name={Reguläre Grammatik},
description={Reguläre Grammatiken sind identisch mit den Typ-3-Grammatiken der Chomsky-Hierarchie. Es sind kontextfreie Grammatiken, die weiteren Einschränkungen unterliegen. Die linke Seite der Regel darf nur ein nichtterminales Symbol beinhalten, die Rechte Seite höchstens ein Nichtterminales.}
}

\newglossaryentry{Semantische Luecke}{%
name={Semantische Lücke},
description={Damit wird der Sachverhalt bezeichnet, dass eine maschinelle oder logische Beschreibung einer Sache nicht alle Aspekte abdeckt, die für einen Nutzer oder das reale Leben relevant sind}
}

\newglossaryentry{Konkrete Syntax}{%
name={Konkrete Syntax},
description={Legt durch Notation (Schlüsselworte, Symbole, \ldots) fest, wie etwas formuliert werden kann. Dazu können reguläre Ausdrücke oder EBNF-Grammatiken verwendet werden}
}

\newglossaryentry{Abstrakte Syntax}{%
name={Abstrakte Syntax},
description={Beschreibt die Struktur der Sprache, bzw. was mit ihr formuliert werden kann. Kann durch abstrakten Syntax-Baum dargestellt werden}
}

\newglossaryentry{Statische Semantik}{%
name={Statische Semantik},
description={Zusätzliche Einschränkungen für die Wohlgeformtheit von Formulierungen. Konsistenzregeln für abstrakte Syntaxgrafen, die zur übersetzungszeit geprüft werden könne. Zum Beispiel müssen alle Bezeichner in einem Bereich eindeutig sein}
}

\newglossaryentry{Dynamische Semantik}{%
name={Dynamische Semantik},
description={Regelt die Bedeutung von Formulierungen, und welches Verhalten von den Formulierungen erzeugt wird. Zum Beispiel wird aus einer Variablendefinition durch den Compiler eine Speicherreservierung}
}

\newglossaryentry{Tagging}{%
name={Tagging},
description={Das Versehen eines Objektes mit einem aussagekräftigen Wort}
}

\newglossaryentry{Tag}{%
name={Tag},
description={Ein aussagekräftiges Wort, das als Meta-Information einem Objekt zugeordnet wird}
}

\usepackage{pgfplots}				           % Package for plots, graphs and diagramms
\pgfplotsset{compat=1.8,width=0.65\columnwidth}


%-------------------- LISTING CONFIGURATION  --------------------%
\definecolor{grey}		{rgb}{0.25, 0.25, 0.25}
\definecolor{lightbeige}{rgb}{0.80, 0.60, 0.30}
\definecolor{lightgrey}	{rgb}{0.90, 0.90, 0.90}
\definecolor{lightgreen}{rgb}{0.00, 0.50, 0.00}
\definecolor{darkgrey}	{rgb}{0.20, 0.20, 0.20}

% Taken from Lena Herrmann at 
% http://lenaherrmann.net/2010/05/20/javascript-syntax-highlighting-in-the-latex-listings-package
\lstdefinelanguage{JavaScript}{
  keywords={typeof, new, true, false, catch, function, return, null, catch, switch, var, if, in, while, do, else, case, break},
  keywordstyle=\color{blue}\bfseries,
  ndkeywords={class, export, boolean, throw, implements, import, this},
  ndkeywordstyle=\color{blue}\bfseries,
  identifierstyle=\color{black},
  sensitive=false,
  comment=[l]{//},
  morecomment=[s]{/*}{*/},
  commentstyle=\color{lightgreen}\ttfamily,
  stringstyle=\color{lightbeige}\ttfamily,
  morestring=[b]',
  morestring=[b]"
}
\lstset{%
	aboveskip			= 3mm,
	belowskip			= 3mm,
	backgroundcolor		= \color{lightgrey},
	basicstyle			= {\scriptsize\ttfamily},
	breakatwhitespace	= true,
	breaklines			= true,
	columns				= flexible,
	commentstyle		= {\color{lightgreen}},
	emptylines			= *1,    % Print >1 empty lines as one (line numbers from origial file)
	frame				= tb,
	identifierstyle		= {\color{darkgrey}},
	inputencoding		= utf8,
	keywordstyle		= {\color{blue}},
	language			= Java,
	numbers				= none,
	numberstyle			= {\tiny\color{grey}},
	showstringspaces	= false,
	stepnumber			= 1,
	stringstyle			= {\color{lightbeige}},
	tabsize				= 4
}


%---------------- Hyperlink- and PDF-Metadata-generation (for PDF/A)
\hypersetup{%
	colorlinks			= false,    % boxed links
	hypertexnames		= false,    % Guessable names for links
	pdftitle			= {\topic},
	pdfauthor			= {\authorName},
	pdfsubject			= {\Subject},
	pdfkeywords			= {\keywords},
	pdfcaptionwriter	= {\authorName},
	pdfcontactaddress	= {\authorStreet},
	pdfcontactcity		= {\authorCity},
	pdfcontactpostcode	= {\authorCityNumber},
	pdfcontactcountry	= {\authorCountry},
	pdfcontactregion	= {\authorFederalState},
	pdfcontactemail		= {\autorEmail},
	pdflang				= {\documentLanguage},
	pdfmetalang			= {\documentLanguage},
	pdfstartview		= {FitV},   % Fit to screen width when opening
	pdftoolbar			= false,    % Hide Toolbar
	pdfdisplaydoctitle	= true,     % Display TOC
	plainpages			= false,    % Helps with links to pages prior to 1
	unicode				= true,     % Allow Unicode characters in bookmarks
}


%---------------- Include Color Profile (from MikTeX)
\immediate\pdfobj stream attr{/N 3} file{\main/sRGB_IEC61966-2-1_black_scaled.icc}
\pdfcatalog{%
/OutputIntents [ <<
/Type /OutputIntent
/S/GTS_PDFA1
/DestOutputProfile \the\pdflastobj\space 0 R
/OutputConditionIdentifier (sRGB IEC61966-2.1)
/Info(sRGB IEC61966-2.1)
>> ]
}



%-----------------------      BEGIN       -----------------------%
\begin{document}

\hyphenation{platt-form-un-ab-hän-gig Re-ti-na Ten-sor-flow}


%-----------------------     Titelseiten
\begin{titlepage}

\includegraphics[height=1.5cm]{HTWG_IN_Modul_Zeichen_pos_2C_rgb}\includegraphics[height=1.5cm]{HTWG_IN_Modul_Zusatz_pos_1C_rgb}\hspace{\fill}

\vspace{2.5cm}

\begin{center}
	\huge{%
		\textbf{\topic} \\[5cm]
	}
	\Large{%
		\textbf{\author}} \\[5cm]
	\large{%
		\textbf{Konstanz, \closingdate} \\[2.3cm]
	}
	\Huge{%
		\textbf{{\textsf{Masterarbeit}}}
	}
\end{center}

\end{titlepage}

\thispagestyle{empty}{%
	\setlength{\parskip}{0.5cm}
	\begin{center}
		\textbf{\huge MASTERARBEIT}
		
		\textbf{zur Erlangung des akademischen Grades}
		
		\textbf{\Large Master of Science (M. Sc.)}
		
		\textbf{an der}
		
		\textsf{\huge Hochschule Konstanz}\\
		{\small Technik, Wirtschaft und Gestaltung}
		
		\textsf{\Large \universityFaculty} \\
		Studiengang \authorStudyTopic%
	\end{center}
}
\begin{center}
	
	\vspace*{2cm}
	
	\begin{tabulary}{\textwidth}{LL}
		\textbf{Thema:} 		& \textbf{\large \topic}\\[5ex]
		\textbf{Verfasser:}		& \authorName\\ 
								& \authorStreet\\ 
								& \authorCityNumber~\authorCity\\[5ex]
		\textbf{1. Prüfer:} 	& \proffessor\\
								& \universityName\\
								& \universityStreet\\
								& \universityCityNumber~\universityCity\\[5ex]
		\textbf{2. Prüfer:}		& \carer\\
								& \companyName\\
								& \companyCityStreet\\
								& \companyCityNumber~\companyCity\\[5ex]
		\textbf{Ausgabe-datum:} 	& \issuedate\\
		\textbf{Abgabe-datum:} 	& \closingdate\\
	\end{tabulary}
\end{center}

\chapter*{Danksagung}
\addcontentsline{toc}{chapter}{Danksagung}

An dieser Stelle, möchte ich mich bei der \universityName\ für die fachliche und kompetente  Betreuung zur Erstellung dieser Abschlussarbeit bedanken, die mir während der Erstellung dieser Arbeit mit Ihrem Wissen zur Seite standen und immer offen gegenüber meiner Fragen waren. Danken möchte ich außerdem allen hier nicht namentlich erwähnten Personen, die mich bei dieser Arbeit auf verschiedenste Weise unterstützt haben.
\vspace*{3cm}
\noindent
Konstanz, \closingdate\hfill \begin{tabular}{c}  \textbf{\authorName} \end{tabular}

\chapter*{Ehrenw\"{o}rtliche Erkl\"{a}rung}
\addcontentsline{toc}{chapter}{Ehrenw\"{o}rtliche Erkl\"{a}rung}

Hiermit erkläre ich,
\textit{\authorName, geboren am \authorBirthday\ in \authorBirthdayCity},\\
\noindent
\begin{tabular}{p{.05\textwidth}p{.8\textwidth}}
	(1) & dass ich meine Bachelorarbeit mit dem Titel \\[1em]
		& \textbf{\topic} \\[1em]
		& bei der \companyName\ unter Anleitung von \proffessor\ selbständig und ohne fremde Hilfe angefertigt habe und keine anderen als die angeführten Hilfen benutzt habe;\\[1em]
	(2) & dass ich die Übernahme wörtlicher Zitate, von Tabellen, Zeichnungen, Bildern und Programmen aus der Literatur oder anderen Quellen (Internet) sowie die Verwendung der Gedanken anderer Autoren an den entsprechenden Stellen innerhalb der Arbeit gekennzeichnet habe.\\[1em]
	(3) & dass die eingereichten Abgabe-Exemplare in Papierform und im PDF-Format vollständig übereinstimmen.\\
\end{tabular}
\vspace*{1cm}

\noindent
Ich bin mir bewusst, dass eine falsche Erklärung rechtliche Folgen haben wird.
\vspace*{3cm}

\noindent
Konstanz, \closingdate\hfill \begin{tabular}{c} \\ \\ \rule{5cm}{1pt} \\ (Unterschrift)
\end{tabular}

\chapter*{Abkürzungsverzeichnis}
\addcontentsline{toc}{chapter}{Abkürzungsverzeichnis}

\begin{acronym}[ANTLR]	% Pass longest acronym as optional parameter to set width of table
\acro{ANTLR}{Another Tool for Language Recognition}
\acro{ADL}{Architecture Description Language}
\acro{API}{Application Programming Interface}
\acro{AST}{Abstract Syntax Tree}
\acro{ATN}{Augmented Transition Network}
\acro{BNF}{Backus-Naur Form}
\acro{CFG}{Context Free Grammar}
\acro{CSV}{Comma Separated Values}
\acro{DDL}{Description Definition Language}
\acro{DSL}{Domain Specific Language}
\acro{EBNF}{Extended \ac{BNF}}
\acro{GPL}{General-purpose Programming Language}
\acro{IDE}{Integrated Development Environment}
\acro{JVM}{Java Virtual Machine}
\acro{LL(k)}{Left to right, left reduction, k symbol lookahead}
\acro{LR(k)}{Left to right, right reduction, k symbol lookahead}
\acro{MOF}{Meta-Object Facility}
\acro{YACC}{Yet Another Compiler Compiler}
\acro{W3C}{World Wide Web Consortium}
\end{acronym}


\tableofcontents\newpage


%-----------------------     Hauptteil
\subfile{01_Main_Part/00_Abstract}
\subfile{01_Main_Part/01_Einleitung}
\subfile{01_Main_Part/02_Einfuehrung}
\subfile{01_Main_Part/03_Main-Part}


%-----------------------     Ende
\subfile{02_Result/00_Conclusion}


%-----------------------     Verzeichnisse
\listoffigures\addcontentsline{toc}{chapter}{Abbildungsverzeichnis}
\listoftables\addcontentsline{toc}{chapter}{Tabellenverzeichnis}
\lstlistoflistings\addcontentsline{toc}{chapter}{Listingverzeichnis}


%\glsaddall%Print all glossary entries
\printglossaries%

\clearpage
\printbibliography\addcontentsline{toc}{chapter}{Literaturverzeichnis}


\end{document}
%-----------------------       END        -----------------------%
